\chapter{Algebraic Structures}

\section{LAWS OF COMPOSITION; ASSOCIATIVITY; COMMUTATIVITY}

\subsection{LAWS OF COMPOSITION}

DEFINITION 1. Let E be a set. A mapping f of E × E into E is called a law of composition on E. The value f(x, y) of f for an ordered pair (x, y) \ensuremath{\in} E × E is called the composition of x and y under this law. A set with a law of composition is called a magma.

The composition of x and y is usually denoted by writing x and y in a definite order and separating them by a characteristic symbol of the law in question (a symbol which it may be agreed to omit). Among the symbols most often used are + and ., the usual convection being to omit the latter if desired; with these symbols the composition of x and y is written respectively as \(x + y\) and x.y or xy. A law denoted by the symbol + is usually called addition (the composition \(x + y\) being called the sum of x and y) and we say that it is written additively; a law denoted by the symbol . is usually called multiplication (the composition x.y = xy being called the product of x and y) and we say that it is written multiplicatively. In the general arguments of paragraphs 1 to 3 of this chapter we shall generally use the symbols \(\top\) and \(\bot\) to denote arbitrary laws of composition.

By an abuse of language, a mapping of a subset of \(E \times E\) into E is sometimes called a law of composition not everywhere defined on E.

Examples. (1) The mappings \((\mathbf{X},\mathbf{Y})\mapsto \mathbf{X}\cup \mathbf{Y}\) and \((\mathbf{X},\mathbf{Y})\mapsto \mathbf{X}\cap \mathbf{Y}\) are laws of composition on the set of subsets of a set E.

\begin{enumerate}
\def\labelenumi{(\arabic{enumi})}
\setcounter{enumi}{1}
\item
  On the set \(\mathbf{N}\) of natural numbers, addition, multiplication and exponentiation are laws of composition (the compositions of \(x \in \mathbf{N}\) and \(y \in \mathbf{N}\) under these laws being denoted respectively by \(x + y\) , \(xy\) or \(x.y\) and \(x^y\) ) (Set Theory, III, § 3, no. 4).
\item
  Let E be a set; the mapping \((\mathbf{X}, \mathbf{Y}) \mapsto \mathbf{X} \circ \mathbf{Y}\) is a law of composition on the set of subsets of \(E \times E\) (Set Theory, II, § 3, no. 3, Definition 6); the mapping \((f, g) \mapsto f \circ g\) is a law of composition on the set of mappings of E into E (Set Theory, II, § 5, no. 2).
\item
  Let E be a lattice (Set Theory, III, § 1, no. 11); if \(\sup(x, y)\) denotes the least upper bound of the set \(\{x, y\}\) , the mapping \((x, y) \mapsto \sup(x, y)\) is a law of composition on E. Similarly for the greatest lower bound \(\inf(x, y)\) . Example 1 above is a particular case of this with \(\mathfrak{P}(E)\) ordered by inclusion.
\item
  Let \((\mathbf{E}_i)_{i\in I}\) be a family of magmas. Let \(\mathsf{T}_i\) denote the law of composition on \(\mathbf{E}_i\) . The mapping
\end{enumerate}

\[
((x _ {i}), (y _ {i})) \mapsto ((x _ {i} \top_ {i} y _ {i}))
\]

is a law of composition on the product \(E = \prod_{i \in I} E_i\) , called the product of the laws \(T_i\) . The set E with this law is called the product magma of the magmas \(E_i\) . In particular, if all the magmas \(E_i\) are equal to the same magma M, we obtain the magma of mappings of I from M.

Let \((x,y)\mapsto x\top y\) be a law of composition on a set E. Given any two subsets X, Y of E, X \(\top\) Y (provided this notation does not lead to confusion†) will denote the set of elements \(x\top y\) in E, where \(x\in X\) , \(y\in Y\) (in other words, the image of X \(\times\) Y under the mapping \((x,y)\mapsto x\top y\) ).

If \(a \in \mathbf{E}\) we usually write \(a \top \mathbf{Y}\) instead of \(\{a\} \top \mathbf{Y}\) and \(\mathbf{X} \top a\) instead of \(\mathbf{X} \top \{a\}\) . The mapping \((\mathbf{X}, \mathbf{Y}) \mapsto \mathbf{X} \top \mathbf{Y}\) is a law of composition on the set of subsets of \(\mathbf{E}\) .

DEFINITION 2. Let E be a magma and \(\top\) denote its law of composition. The law of composition \((x, y) \mapsto y \top x\) on E is called the opposite of the above. The set E with this law is called the opposite magma of E.

Let E and E' be two magmas; we shall denote their laws by the same symbol T. Conforming with the general definitions (Set Theory, IV, § 1, no. 5), a bijective mapping f of E onto E' such that

\[
f (x \top y) = f (x) \top f (y)\tag{1}
\]

for every ordered pair \((x, y) \in \mathbf{E} \times \mathbf{E}\) is called an isomorphism of \(\mathbf{E}\) onto \(\mathbf{E}'\) . \(\mathbf{E}\) and \(\mathbf{E}'\) are said to be isomorphic if there exists an isomorphism of \(\mathbf{E}\) onto \(\mathbf{E}'\) .

More generally:

DEFINITION 3. A mapping \(f\) of \(\mathbf{E}\) into \(\mathbf{E}'\) such that relation (1) holds for every ordered pair \((x, y) \in \mathbf{E} \times \mathbf{E}\) is called a homomorphism, or morphism, of \(\mathbf{E}\) into \(\mathbf{E}'\) ; if \(\mathbf{E} = \mathbf{E}'\) , \(f\) is called an endomorphism of \(\mathbf{E}\) .

The identity mapping of a magma E is a homomorphism, the composition of two homomorphism is a homomorphism.

For a mapping \(f\) of \(E\) into \(E'\) to be an isomorphism, it is necessary and sufficient that it be a bijective homomorphism and \(\overline{f}^{-1}\) is then an isomorphism of \(E'\) onto \(E\) .

\subsection{COMPOSITION OF AN ORDERED SEQUENCE OF ELEMENTS}

Recall that a family of elements of a set E is a mapping \(t \mapsto x_{t}\) of a set I (called an indexing set) into E; a family \((x_{t})_{t \in I}\) is called finite if the indexing set is finite.

A finite family \((x_{t})_{t\in I}\) of elements of E whose indexing set I is totally ordered is called an ordered sequence of elements of E.

In particular, every finite sequence \((x_{i})_{i\in\mathbb{H}}\) , where H is a finite subset of the set N of natural numbers, can be considered as an ordered sequence if H is given the order relation induced by the relation \(m \leqslant n\) between natural numbers.

Two ordered sequences \((x_{i})_{i\in\mathbf{I}}\) and \((y_{k})_{k\in\mathbf{K}}\) are called similar if there exists an ordered set isomorphism \(\phi\) of I onto K such that \(y_{\phi(i)} = x_{i}\) for all \(i \in I\) .

Every ordered sequence \((x_{\alpha})_{\alpha \in \mathbf{A}}\) is similar to a suitable finite sequence. For there exists an increasing bijection of \(\mathbf{A}\) onto an interval \([0, n]\) of \(\mathbf{N}\) .

DEFINITION 4. Let \((x_{\alpha})_{\alpha \in \mathbf{A}}\) be an ordered sequence of elements in a magma \(\mathbf{E}\) whose indexing set \(\mathbf{A}\) is non-empty. The composition (under the law \(\top\) ) of the ordered sequence \((x_{\alpha})_{\alpha \in \mathbf{A}}\) , denoted by \(\bigcap_{\alpha \in \mathbf{A}} x_{\alpha}\) , is the element of \(\mathbf{E}\) defined by induction on the number of elements in \(\mathbf{A}\) as follows:

\begin{enumerate}
\def\labelenumi{(\arabic{enumi})}
\item
  if \(\mathbf{A} = \{\beta\}\) then \(\bigcap_{\alpha \in \mathbf{A}} x_{\alpha} = x_{\beta}\) ;
\item
  if A has \(p > 1\) elements, \(\beta\) is the least element of A and \(A' = A - \{\beta\}\) , then \(\bigcap_{\alpha \in A} x_{\alpha} = x_{\beta} \top \left( \bigcap_{\alpha \in A} x_{\alpha} \right)\) .
\end{enumerate}

It follows immediately (by induction on the number of elements in the indexing sets) that the compositions of two similar ordered sequences are equal; in particular, the composition of any ordered sequence is equal to the composition of a finite sequence. If \(\mathbf{A} = \{\lambda, \mu, \nu\}\) has three elements ( \(\lambda < \mu < \nu\) ) the composition \(\bigcap_{\alpha \in \mathbf{A}} x_{\alpha}\) is \(x_{\lambda} \top (x_{\mu} \top x_{\nu})\) .

Remark. Note that there is a certain arbitrariness about the definition of the composition of an ordered sequence; the induction we introduced proceeds ``from right to left''. If we proceeded ``from left to right'', the composition of the above ordered sequence \((x_{\lambda}, x_{\mu}, x_{\nu})\) would be \((x_{\lambda} \top x_{\mu}) \top x_{\nu}\) .

\[
\left(x _ {\alpha}\right) _ {\alpha \in A}
\]

\[
\begin{array}{c} \underline {{\mathsf {I}}} \\ \alpha \in \mathbf {A} \end{array}
\]

denoted by \(\sum_{\alpha \in \mathbf{A}} x_{\alpha}\) and called the sum of the ordered sequence \((x_{\alpha})_{\alpha \in \mathbf{A}}\) (the \(x_{\alpha}\) being called the terms of the sum); for a law written multiplicatively it is usually denoted by \(\prod_{\alpha \in A} x_{\alpha}\) and called the product of the ordered sequence \((x_{\alpha})\) (the \(x_{\alpha}\) being called the factors of the product).†

When there is no possible confusion over the indexing set (nor over its ordering) it is often dispensed with in the notation for the composition of an ordered sequence and we then write, for example for a law written additively, \(\sum_{\alpha} x_{\alpha}\) instead of \(\sum_{\alpha \in A} x_{\alpha}\) ; similarly for the other notations.

For a law denoted by \(\top\) the composition of a sequence \((x_{i})\) with indexing set a non-empty interval \((p,q)\) of \(\mathbf{N}\) is denoted by \(\bigcap_{p\leqslant i\leqslant q}x_{i}\) or \(\bigcap_{i = p}^{q}x_{i}\) ; similarly for laws denoted by other symbols.

Let E and F be two magmas whose laws of composition are denoted by T and f a homomorphism of E into F. For every ordered sequence \((x_{\alpha})_{\alpha \in \mathbf{A}}\) of elements of E

\[
f \left(\underset {\alpha \in A} {T}\right) = \underset {\alpha \in A} {T} f (x _ {\alpha}).\tag{2}
\]

\subsection{ASSOCIATIVE LAWS}

DEFINITION 5. A law of composition \((x,y)\mapsto x\top y\) on a set \(\mathbf{E}\) is called associative if, for all elements \(x,y,z\) in \(\mathbf{E}\) ,

\[
(x \top y) \top z = x \top (y \top z).
\]

A magma whose law is associative is called an associative magma.

The opposite law of an associative law is associative.

Examples. (1) Addition and multiplication of natural numbers are associative laws of composition on \(\mathbf{N}\) (Set Theory, III, § 3, no. 3, Corollary to Proposition 5)

\begin{enumerate}
\def\labelenumi{(\arabic{enumi})}
\setcounter{enumi}{1}
\tightlist
\item
  The laws cited in Examples (1), (3) and (4) of no. 1 are associative.
\end{enumerate}

THEOREM 1 (Associativity theorem). Let E be an associative magma whose law is denoted by \(\top\) . Let A be a totally ordered non-empty finite set, which is the union of an ordered sequence of non-empty subsets \((\mathbf{B}_l)_{i \in I}\) such that the relations \(\alpha \in \mathbf{B}_i\) , \(\beta \in \mathbf{B}_j\) , \(i < j\) imply \(\alpha < \beta\) ; let \((x_\alpha)_{\alpha \in A}\) be an ordered sequence of elements in E with A as indexing set. Then

\[
\bigcap_ {\alpha \in A} x _ {\alpha} = \bigcap_ {i \in I} \left(\bigcap_ {\alpha \in B _ {i}} x _ {\alpha}\right)\tag{3}
\]

† The use of this terminology and the notation \(\prod_{\alpha \in \mathbf{A}} x_{\alpha}\) must be avoided if there is any risk of confusion with the product of the sets \(x_{\alpha}\) defined in the theory of sets (Set Theory, II, § 5, no. 3). However, if the \(x_{\alpha}\) are cardinals and addition (resp. multiplication) is the cardinal sum (resp. the cardinal product), the cardinal denoted by \(\sum_{\alpha \in \mathbf{A}} x_{\alpha}\) (resp. \(\prod_{\alpha \in \mathbf{A}} x_{\alpha}\) ) in the above notation is the cardinal sum (resp. cardinal product) of the family \((x_{\alpha})_{\alpha \in \mathbf{A}}\) (Set Theory, III, § 3, no. 3).

We prove the theorem by induction on the cardinal n of A. Let p be the cardinal of I and h its least element; let \(J = I - \{h\}\) . If n = 1, as the \(B_{i}\) are non-empty, of necessity p = 1 and the theorem is obvious. Otherwise, assuming the theorem holds for an indexing set with at most n - 1 elements, we distinguish two cases:

\begin{enumerate}
\def\labelenumi{(\alph{enumi})}
\tightlist
\item
  \(B_{h}\) has a single element \(\beta\) . Let \(C = \bigcup_{i \in J} B_{i}\) . The left-hand side of (3) is equal, by definition, to \(x_{\beta} \top \left( \bigcap_{\alpha \in C} x_{\alpha} \right)\) ; the right-hand side is equal, by definition, to
\end{enumerate}

\[
x _ {\beta} \top \left(\underset {i \in J} {\mathsf {T}} \left(\underset {\alpha \in B _ {i}} {\mathsf {T}} x _ {\alpha}\right)\right);
\]

the result follows from the fact that the theorem is assumed true for C and \((\mathbf{B}_{i})_{i\in J}\) .

\begin{enumerate}
\def\labelenumi{(\alph{enumi})}
\setcounter{enumi}{1}
\tightlist
\item
  Otherwise, let \(\beta\) be the least element of A (and hence of \(\mathbf{B}_h\) ); let \(\mathbf{A}' = \mathbf{A} - \{\beta\}\) and let \(\mathbf{B}_i' = \mathbf{A}' \cap \mathbf{B}_i\) for \(i \in \mathbf{I}\) ; then \(\mathbf{B}_i' = \mathbf{B}_i\) for \(i \in \mathbf{J}\) . The set \(\mathbf{A}'\) has \(n - 1\) elements and the conditions of the theorem hold for \(\mathbf{A}'\) and its subsets \(\mathbf{B}_i'\) ; hence by hypothesis:
\end{enumerate}

\[
\underset {\alpha \in A ^ {\prime}} {T} x _ {\alpha} = \left(\underset {\alpha \in B ^ {\prime} h} {T} x _ {\alpha}\right) \top \left(\underset {i \in J} {T} \left(\underset {\alpha \in B _ {i}} {T} x _ {\alpha}\right)\right).
\]

Forming the composition of \(x_{\beta}\) with each side, we have on the left-hand side side, by definition, \(\bigcap_{\alpha \in A} x_{\alpha}\) and on the right-hand side, using the associativity,

\[
\left(x _ {\beta} \top \left(\underset {\alpha \in B ^ {\prime} h} {\top} x _ {\alpha}\right)\right) \top \left(\underset {i \in J} {\top} \left(\underset {\alpha \in B _ {i}} {\top} x _ {\alpha}\right)\right)
\]

which is equal, by Definition 3, to the right-hand side of formula (3).

For an associative law denoted by \(\top\) the composition \(\bigcap_{p\leqslant i\leqslant q}x_i\) of a sequence \((x_i)_{i\in [p,q]}\) is also denoted (since no confusion can arise) by

\[
x _ {p} \top \dots \top x _ {q}.
\]

A particular case of Theorem 1 is the formula

\[
x _ {0} \top x _ {1} \top \dots \top x _ {n} = (x _ {0} \top x _ {1} \top \dots \top x _ {n - 1}) \top x _ {n}.
\]

Consider an ordered sequence of n terms all of whose terms are equal to the same element \(x \in E\) . The composition of this sequence is denoted by \(\top^{n} x\) for a law denoted by \(\top, \bot^{n} x\) for a law denoted by \(\bot\) . For a law written multiplicatively the composition is denoted by \(x^{n}\) and called the n-th power of x. For a law written additively the composition is usually denoted by nx. The associativity theorem applied to an ordered sequence all of whose terms are equal gives the equation

\[
\begin{array}{c} n _ {1} + n _ {2} + \dots + n _ {p} \\ \top \end{array} x = \binom {n _ {1}} {\top} x) \top \binom {n _ {2}} {\top} x) \top \dots \top \binom {n _ {p}} {\top} x).
\]

In particular, if \(p = 2\) ,

\[
{ } ^ { m + n } \top x = \left( \begin{array} { c } m \\ \top x \end{array} \right) \top \left( \begin{array} { c } n \\ \top x \end{array} \right)\tag{4}
\]

and if \(n_1 = n_2 = \cdots = n_p = m,\)

\[
\frac {p m}{T} x = \frac {p}{T} \left(\frac {m}{T} x\right).\tag{5}
\]

If X is a subset of E, we sometimes denote, in conformity with the above notation, by \(\overset{p}{\top}X\) the set \(X_{1}\top X_{2}\top\cdots\top X_{p}\) , where

\[
\mathbf {X} _ {1} = \mathbf {X} _ {2} = \dots = \mathbf {X} _ {p} = \mathbf {X};
\]

it is thus the set of all compositions \(x_{1} \top x_{2} \top \cdots \top x_{p}\) with \(x_{1} \in X, x_{2} \in X, \ldots, x_{p} \in X\) .

It is important not to confuse this set with the set of \(\top x\) , where x runs through X.

\subsection{STABLE SUBSETS. INDUCED LAWS}

DEFINITION 6. A subset A of a set E is called stable under the law of composition \(\top\) on E if the composition of two elements of A belongs to A. The mapping \((x,y)\mapsto x\top y\) of \(\mathbf{A}\times \mathbf{A}\) into A is then called the law induced on A by the law \(\top\) . The set A with the law induced by \(\top\) is called a submagma of E.

In other words, for A to be stable under a law \(\top\) it is necessary and sufficient that \(A \top A \subset A\) . A stable subset of E and the corresponding submagma are often identified.

The intersection of a family of stable subsets of E is stable; in particular there exists a smallest stable subset A of E containing a given subset X; it is said to be generated by X and X is called a generating system of A or a generating set of A. The corresponding submagma is also said to be generated by X.

PROPOSITION 1. Let E and F be two magmas and f a homomorphism of E into F.

\begin{enumerate}
\def\labelenumi{(\roman{enumi})}
\item
  The image under \(f\) of a stable subset of \(\mathbf{E}\) is a stable subset of \(\mathbf{F}\) .
\item
  The inverse image under \(f\) of a stable subset \(\mathbf{F}\) is a stable subset of \(\mathbf{E}\) .
\item
  Let \(\mathbf{X}\) be a subset of \(\mathbf{E}\) . The image under \(f\) of the stable subset of \(\mathbf{E}\) generated by \(\mathbf{X}\) is the stable subset of \(\mathbf{F}\) generated by \(f(\mathbf{X})\) .
\item
  If \(g\) is a second homomorphism of \(\mathbf{E}\) from \(\mathbf{F}\) the set of elements \(x\) of \(\mathbf{E}\) such that \(f(x) = g(x)\) is a stable subset of \(\mathbf{E}\) .
\end{enumerate}

\subsection{PERMUTABLE ELEMENTS. COMMUTATIVE LAWS}

\begin{enumerate}
\def\labelenumi{(\arabic{enumi})}
\setcounter{enumi}{3}
\tightlist
\item
  Let \((x, y) \mapsto x \top y\) be a commutative law on E; the law
\end{enumerate}

\[
(\mathbf {X}, \mathbf {Y}) \mapsto \mathbf {X} \top \mathbf {Y}
\]

between subsets of E is commutative.

DEFINITION 9. Let E be a magma and X a subset of E. The set of elements of E which commute with each of the elements of X is called the centralizer of X.

Let X and Y be two subsets of E and X' and Y' their respective centralizers. If X ⊂ Y, then Y' ⊂ X'.

Let \((\mathbf{X}_i)_{i\in I}\) be a family of subsets of \(E\) and for all \(i\in I\) let \(\mathbf{X}_i'\) be the centralizer of \(\mathbf{X}_i\) . The centralizer of \(\bigcup_{i\in I}\mathbf{X}_i\) is \(\bigcap_{i\in I}\mathbf{X}_i'\) .

Let X be a subset of E and \(X'\) the centralizer of X. The centralizer \(X''\) of \(X'\) is called the bicentralizer of X. Then \(X \subset X''\) . The centralizer \(X'''\) of \(X''\) is equal to \(X'\) . For \(X'\) is contained in its bicentralizer \(X'''\) and the relation \(X \subset X''\) implies \(X''' \subset X'\) .

PROPOSITION 3. Let E be an associative magma whose law is denoted by \(\top\) . If an element \(x\) of E commutes with each of the elements \(y\) and \(z\) of E, it commutes with \(y \top z\) .

For

\[
\begin{array}{r l} x \top (y \top z) = (x \top y) \top z & = (y \top x) \top z \\ & = y \top (x \top z) = y \top (z \top x) = (y \top z) \top x. \end{array}
\]

COROLLARY. Let E be an associative magma. The centralizer of any subset of E is a stable subset of E.

DEFINITION 10. The centralizer of a magma E is called the centre of E. An element of the centre of E is called a central element of E.

If E is an associative magma its centre is a stable subset by the Corollary to Proposition 3 and the law induced on its centre is commutative.

PROPOSITION 4. Let E be an associative magma, X and Y two subsets of E. If every element of X commutes with every element of Y every element of the stable subset generated by X commutes with every element of the stable subset generated by Y.

Let \(X'\) and \(X''\) be the centralizer and bicentralizer of X. They are stable subsets of E. Now \(X \subset X''\) and \(Y \subset X'\) and hence \(X''\) (resp. \(X'\) ) contains the stable subset of E generated by X (resp. Y). As every element of \(X''\) commutes with every element of \(X'\) , the proposition follows.

COROLLARY 1. If \(x\) and \(y\) are permutable under the associative law \(\top\) so are \(\top^m x\) and \(\top^n y\) for all integers \(m > 0\) and \(n > 0\) ; in particular \(\top^m x\) and \(\top^n x\) are permutable for all \(x\) and all integers \(m > 0\) and \(n > 0\) .

COROLLARY 2. If all pairs of elements of a subset X are permutable under an associative law T, the law induced by T on the stable subset generated by X is associative and commutative.

THEOREM 2 (Commutativity theorem). Let \(\top\) be an associative law of composition on E; let \((x_{\alpha})_{\alpha \in A}\) be a non-empty finite family of elements of E which are pairwise permutable; let B and C be two totally ordered sets with A as underlying set. Then \(\bigcap_{\alpha \in B} x_{\alpha} = \bigcap_{\alpha \in C} x_{\alpha}\) .

Since the theorem is true if A has a single element, we argue by induction on the number p of elements in A. Let p be an integer \textgreater1 and suppose the theorem is true when Card A \textless{} p.~We prove it for Card A = p.~It may be assumed that A is the interval \([0, p - 1]\) in N; the composition of the ordered sequence \((x_{\alpha})_{\alpha \in \mathbf{A}}\) defined by the natural order relation on \(\mathbf{A}\) is \(\prod_{i = 0}^{p - 1}x_i\) .

Let A be given another total ordering and let h be the least element of A under this ordering and \(A'\) the set of other elements of A (totally ordered by the induced ordering). Suppose first 0 \textless{} h \textless{} p - 1 and let \(P = \{0, 1, \ldots, h - 1\}\) and \(Q = \{h + 1, \ldots, p - 1\}\) ; the theorem being assumed true for \(A'\) , applying the associativity theorem, we obtain (since \(A' = P \cup Q\) )

\[
\bigcap_ {\alpha \in \mathbf {A} ^ {\prime}} x _ {\alpha} = \left(\prod_ {i = 0} ^ {h - 1} x _ {i}\right) \top \left(\prod_ {i = h + 1} ^ {p - 1} x _ {i}\right)
\]

whence, composing \(x_{h}\) with both sides and repeatedly applying the commutativity and associativity of T:

\[
\begin{array}{r l} \underset {\alpha \in \mathbf {A}} {\top} x _ {\alpha} & = x _ {h} \top \left(\underset {\alpha \in \mathbf {A} ^ {\prime}} {\top} x _ {\alpha}\right) = x _ {h} \top \left(\underset {i = 0} {\overset {h - 1} {\top}} x _ {i}\right) \top \left(\underset {i = h + 1} {\overset {p - 1} {\top}} x _ {i}\right) \\ & = \left(\underset {i = 0} {\overset {h - 1} {\top}} x _ {i}\right) \top x _ {h} \top \left(\underset {i = h + 1} {\overset {p - 1} {\top}} x _ {i}\right) = \underset {i = 0} {\overset {p - 1} {\top}} x _ {i}, \end{array}
\]

which proves the theorem for this case. If h = 0 or h = p - 1, the same result follows, but more simply, the terms arising from P or the terms arising from Q not appearing in the formulae.

Under a commutative associative law on a set E the composition of a finite family \((x_{\alpha})_{\alpha \in A}\) of elements of E is by definition the common value of the composition of all the ordered sequences obtained by totally ordering A in all possible ways. This composition will still be denoted by \(\bigcap_{\alpha \in A} x_{\alpha}\) under a law denoted by T; similarly for other notations.

THEOREM 3. Let \(\top\) be an associative law on E and \((x_{\alpha})_{\alpha \in \mathbf{A}}\) a non-empty finite family of elements of E which are pairwise permutable. If A is a union of non-empty subsets \((\mathbf{B}_i)_{i\in I}\) which are pairwise disjoint, then

\[
\bigcap_ {\alpha \in A} x _ {\alpha} = \bigcap_ {i \in I} \left(\bigcap_ {\alpha \in B _ {i}} x _ {\alpha}\right).\tag{6}
\]

This follows from Theorem 2 if A and I are totally ordered so that the \(B_{i}\) satisfy the conditions of Theorem 1.

We single out two important special cases of this theorem:

\begin{enumerate}
\def\labelenumi{\arabic{enumi}.}
\tightlist
\item
  If \((x_{\alpha\beta})_{(\alpha,\beta)\in\mathbf{A}\times\mathbf{B}}\) is a finite family of permutable elements of an associative magma whose indexing set is the product of two non-empty finite sets A, B (a ``double family''), then
\end{enumerate}

\[
\underset {(\alpha , \beta) \in A \times B} {\top} x _ {\alpha \beta} = \underset {\alpha \in A} {\top} \left(\underset {\beta \in B} {\top} x _ {\alpha \beta}\right) = \underset {\beta \in B} {\top} \left(\underset {\alpha \in A} {\top} x _ {\alpha \beta}\right)\tag{7}
\]

as follows from Theorem 3 by considering \(\mathbf{A} \times \mathbf{B}\) as the union of the sets \(\{\alpha\} \times \mathbf{B}\) on the one hand and of the sets \(\mathbf{A} \times \{\beta\}\) on the other.

For example, if B has n elements and for each \(\alpha \in A\) all the \(x_{\alpha\beta}\) have the same value \(x_{\alpha}\) , then

\[
\bigcap_ {\alpha \in \mathbf {A}} \left(\frac {n}{T} x _ {\alpha}\right) = \frac {n}{T} \left(\bigcap_ {\alpha \in \mathbf {A}} x _ {\alpha}\right).\tag{8}
\]

If B has two elements, we obtain the following results: let \((x_{\alpha})_{\alpha \in A}, (y_{\alpha})_{\alpha \in A}\) be two non-empty families of elements of E. If the \(x_{\alpha}\) and the \(y_{\beta}\) are pairwise permutable, then

\[
\underset {\alpha \in \mathbf {A}} {\top} x _ {\alpha} \top y _ {\alpha} = \left(\underset {\alpha \in \mathbf {A}} {\top} x _ {\alpha}\right) \top \left(\underset {\alpha \in \mathbf {A}} {\top} y _ {\alpha}\right).\tag{9}
\]

Because of formula (7) the composition of a double sequence \((x_{ij})\) whose indexing set is the product of two intervals \(\{p, q\}\) and \(\{r, s\}\) in \(\mathbf{N}\) is often denoted for a commutative associative law written additively by

\[
\sum_ {i = p} ^ {q} \sum_ {j = r} ^ {s} x _ {i j} \quad \text { or } \quad \sum_ {j = r} ^ {s} \sum_ {i = p} ^ {q} x _ {i j}
\]

and similarly for laws denoted by other symbols.

\begin{enumerate}
\def\labelenumi{\arabic{enumi}.}
\setcounter{enumi}{1}
\tightlist
\item
  Let \(n\) be an integer \(>0\) and let \(A\) be the set of ordered pairs of integers \((i,j)\) such that \(0 \leqslant i \leqslant n, 0 \leqslant j \leqslant n\) and \(i < j\) ; the composition of a family \((x_{ij})_{(i,j) \in A}\) (under a commutative associative law) is also denoted by \(\bigcap_{0 \leqslant i < j \leqslant n} x_{ij}\) (or simply \(\bigcap_{i < j} x_{ij}\) , if no confusion arises); Theorem 3 here gives the formulae
\end{enumerate}

\[
\underset {0 \leqslant i <   j \leqslant n} {\mathsf {T}} x _ {i j} = \underset {i = 0} {\overset {n - 1} {\mathsf {T}}} \left(\underset {j = i + 1} {\overset {n} {\mathsf {T}}} x _ {i j}\right) = \underset {j = 1} {\overset {n} {\mathsf {T}}} \left(\underset {i = 0} {\overset {j - 1} {\mathsf {T}}} x _ {i j}\right).\tag{10}
\]

There are analogous formulae to (7) for a family whose indexing set is the product of more than two sets and analogous formulae to (10) for a family whose indexing set is the set \(S_{p}\) of strictly increasing sequences \((i_{k})_{1\leqslant k\leqslant p}\) of p integers such that \(0\leqslant i_{k}\leqslant n\) ( \(p\leqslant n+1\) ): in the latter case the composition of the family \((x_{i_{1}i_{2}\ldots i_{p}})(i_{1},\ldots,i_{p})\in\mathbb{S}_{p}\) is denoted by

\[
\underset {0 \leqslant i _ {1} <   i _ {2} <   \dots <   i _ {p} <   n} {\mathsf {T}} x _ {i _ {1} i _ {2} \ldots i _ {p}}, \quad \text {or simply} \quad \underset {i _ {1} <   i _ {2} <   \dots <   i _ {p}} {\mathsf {T}} x _ {i _ {1} i _ {2} \ldots i _ {p}}.
\]

PROPOSITION 5. Let E and F be magmas whose laws are denoted by \(\top\) and let \(f\) and \(g\) be homomorphisms of E into F. Let \(f \top g\) be the mapping \(x \mapsto f(x) \top g(x)\) of E into F. If F is associative and commutative, \(f \top g\) is a homomorphism.

For all elements \(x\) and \(y\) of \(\mathbf{E}\) :

\[
\begin{array}{r l} (f \top g) (x \top y) & = f (x \top y) \top g (x \top y) = f (x) \top f (y) \top g (x) \top g (y) \\ & = f (x) \top g (x) \top f (y) \top g (y) = ((f \top g) (x)) \top ((f \top g) (y)). \end{array}
\]

\subsection{QUOTIENT LAWS}

DEFINITION 11. Let E be a set. A law of composition \(\top\) and an equivalence relation R on E are said to be compatible if the relations \(x \equiv x'\) (mod R) and \(y \equiv y'\) (mod R) (for \(x, x', y, y' \in \mathbf{E}\) ) imply \(x \top y \equiv x' \top y'\) (mod R); the law of composition on the quotient set E/R which maps the equivalence classes of \(x\) and \(y\) to the equivalence class of \(x \top y\) is called the quotient law of the law \(\top\) with respect to R. The set E/R with the quotient law is called the quotient magma of E with respect to R.

To say that an equivalence relation R on E is compatible with the internal law of composition \(f: E \times E \to E\) on E means that the mapping f is compatible (in the sense of Set Theory, II, §6, no. 5) with the product equivalence relation \(R \times R\) on \(E \times E\) and the equivalence relation R on E. (Set Theory, II, §6, no. 8). This also means that the graph of R is a submagma of \(E \times E\) .

If the law \(\top\) is associative (resp. commutative) so is the quotient law (more briefly we say that associativity, or commutativity, is preserved when passing to the quotient).

The canonical mapping from the magma E to the magma E/R is a homomorphism.

For a mapping g of E/R into a magma F to be a homomorphism it is necessary and sufficient that the composition of g with the canonical mapping of E onto E/R be a homomorphism.

The two following propositions are immediate from the definitions:

PROPOSITION 6. Let E and F be two magmas and f a homomorphism of E into F. Let R\{x,y\} denote the relation \(f(x)=f(y)\) between elements x,y of E. Then R is an equivalence relation on E compatible with the law on E and the mapping of E/R onto f(E) derived from f by passing to the quotient is an isomorphism of the quotient magma E/R onto the submagma f(E) of F.

PROPOSITION 7. Let E be a magma and R an equivalence relation on E compatible with the law on E. For an equivalence relation S on E/R to be compatible with the quotient law it is necessary and sufficient that S be of the form T/R where T is an equivalence relation on E implied by R and compatible with the law on E. The canonical mapping of E/T onto (E/R)/(T/R) (Set Theory, II, §6, no. 7) is then a magma isomorphism.

PROPOSITION 8. Let E be a magma, A a stable subset of E and R an equivalence relation on E compatible with the law on E. The saturation B of A with respect to R (Set

Theory, II, § 6, no. 5) is a stable subset. The equivalence relations \(\mathbf{R}_{\mathbf{A}}\) and \(\mathbf{R}_{\mathbf{B}}\) induced by \(\mathbf{R}\) on \(\mathbf{A}\) and \(\mathbf{B}\) respectively are compatible with the induced laws and the mapping derived from the canonical injection of \(\mathbf{A}\) into \(\mathbf{B}\) by passing to the quotients is a magma isomorphism of \(\mathbf{A} / \mathbf{R}_{\mathbf{A}}\) onto \(\mathbf{B} / \mathbf{R}_{\mathbf{B}}\) .

Let \(\top\) denote the law on E. If \(x\) and \(y\) are two elements of B there exist two elements \(x'\) and \(y'\) of A such that \(x \equiv x' \pmod{R}\) and \(y \equiv y' \pmod{R}\) ; then \(x \top y \equiv x' \top y' \pmod{R}\) and \(x' \top y' \in A\) , whence \(x \top y \in B\) . Thus B is a stable subset of E and the other assertions are obvious.

Let M be a magma and \(((u_{\alpha}, v_{\alpha}))_{\alpha \in I}\) a family of elements of \(M \times M\) . Consider all the equivalence relations S on M which are compatible with the law on M and such that \(u_{\alpha} \equiv v_{\alpha} \pmod{S}\) for all \(\alpha \in I\) . The intersection of the graphs of these relations is the graph of an equivalence relation R which is compatible with the law on M and such that \(u_{\alpha} \equiv v_{\alpha} \pmod{R}\) . Hence R is the finest (Set Theory, III, §1, nos. 3 and 7) equivalence relation with these two properties. It is called the equivalence relation compatible with the law on M generated by the \((u_{\alpha}, v_{\alpha})\) .

PROPOSITION 9. Preserving the above notation, let \(f\) be a homomorphism of \(\mathbf{M}\) into a magma such that \(f(u_{\alpha}) = f(v_{\alpha})\) for all \(\alpha \in \mathbf{I}\) . Then \(f\) is compatible with \(\mathbf{R}\) .

Let T be the equivalence relation associated with f.~Then \(u_{\alpha} \equiv v_{\alpha} \pmod{T}\) for all \(\alpha \in I\) and T is compatible with the law on M, hence T is coarser than R; this proves the proposition.

\section{IDENTITY ELEMENT; CANCELLABLE ELEMENTS; INVERTIBLE ELEMENTS}

\subsection{IDENTITY ELEMENT}

DEFINITION 1. Under a law of composition \(\top\) on a set \(\mathbf{E}\) an element \(e\) of \(\mathbf{E}\) is called an identity element if, for all \(x \in \mathbf{E}\) , \(e \top x = x \top e = x\) .

There exists at most one identity element under a given law \(\top\) , for if \(e\) and \(e'\) are identity elements then \(e = e \top e' = e'\) . An identity element is permutable with every element: it is a central element.

DEFINITION 2. A magma with an identity element is called a unital magma. If E, E' are unital magmas, a homomorphism of the magma E into the magma E' which maps the identity element of E to the identity element of E' is called a unital homomorphism (or morphism) of E into E'. An associative unital magma is called a monoid.

If E, E' are monoids, a unital morphism of E into E' is called a monoid homomorphism or a monoid morphism of E into E'.

Examples. (1) In the set \(\mathbf{N}\) of natural numbers 0 is an identity element under addition and 1 is an identity element under multiplication. Each of these two laws gives \(\mathbf{N}\) a commutative monoid structure (Set Theory, III, §3, no. 3).

\begin{enumerate}
\def\labelenumi{(\arabic{enumi})}
\setcounter{enumi}{1}
\item
  In the set of subsets of a set E, \(\varnothing\) is an identity element under the law \(\cup\) and E under the law \(\cap\) . More generally, in a lattice the least element, if it exists, is identity element under the law sup; conversely, if there exists an identity element under this law it is the least element of the set. Similarly for the greatest element and the law inf.
\item
  The set \(\mathbf{N}\) has no identity element under the law \((x, y) \mapsto x^y\) . Under the law \((\mathbf{X}, \mathbf{Y}) \mapsto \mathbf{X} \circ \mathbf{Y}\) between subsets of \(\mathbf{E} \times \mathbf{E}\) the diagonal \(\Delta\) is the identity element. Under the law \((f, g) \mapsto f \circ g\) between mappings of \(\mathbf{E}\) into \(\mathbf{E}\) the identity mapping of \(\mathbf{E}\) onto \(\mathbf{E}\) is the identity element.
\item
  Let E be a magma and R an equivalence relation on E compatible with the law on E (§ 1, no. 6). If e is an identity element of E the canonical image of e in E/R is an identity element of the magma E/R.
\end{enumerate}

The identity element of a unital magma is a unital homomorphism; the composition of two unital homomorphisms is also one. For a mapping to be a unital magma isomorphism it is necessary and sufficient that it be a bijective unital homomorphism and the inverse mapping is then a unital homomorphism. Let E and E' be unital magmas and \(e'\) the identity element of \(E'\) ; the constant mapping of E into \(E'\) mapping E to \(e'\) is a unital homomorphism, called a trivial homomorphism.

The product of a family of unital magmas (resp. monoids) is a unital magma (resp. monoid).

Every quotient magma of a unital magma (resp. monoid) is a unital magma (resp. monoid).

Let E be a unital magma and e its identity element. A submagma A of E such that \(e \in A\) is called a unital submagma of E. Clearly e is the identity element of the magma A. Every intersection of unital submagmas of E is a unital submagma of E. If X is a subset of E then there exists a smallest unital submagma of E containing X; it is called the unital submagma of E generated by X; it is equal to \(\{e\}\) if X is empty. If E is a monoid, a unital submagma of E is called a submonoid of E.

If F is a magma without identity element a submagma of F may possess an identity element. For example, if F is associative and h is an idempotent element of F (I, § 1, no. 4), the set of \(h \top x \top h\) , where x runs through F, is a submagma of F with h as identity element.

If E is a magma with identity element e, a submagma A of E such that \(e \notin A\) may still possess an identity element.

DEFINITION 3. Let E be a unital magma. The identity element of E is called the composition of the empty family of elements of E.

If \((x_{\alpha})_{\alpha \in \varnothing}\) is the empty family of elements of E its composition \(e\) is also denoted by \(\bigcap_{\alpha \in \varnothing} x_{\alpha}\) . For example, we write

\[
\bigcap_ {q \leqslant i \leqslant p} x _ {i} = e
\]

when \(p < q\) ( \(p, q \in \mathbf{N}\) ). Similarly we write \(\top^0 x = e\) for arbitrary \(x\) . With these definitions Theorems 1 and 3 of § 1 remain true if the hypothesis that the sets \(\mathbf{A}\) and \(\mathbf{B}_i\) are non-empty is suppressed. Similarly the formulae \(\top^{m+n} x = \left( \frac{m}{\top} x \right)^{\top} \left( \frac{n}{\top} x \right)\) and \(\top^{mn} x = \frac{m}{\top} \left( \frac{n}{\top} x \right)\) are then true for \(m \geqslant 0, n \geqslant 0\) .

Let E be a unital magma whose law is denoted by \(\top\) and \(e\) its identity element. The support of a family \((x_{i})_{i\in I}\) of elements of E is the set of indices \(i\in I\) such that \(x_{i}\neq e\) . Let \((x_{i})_{i\in I}\) be a family of elements of E with finite support. We shall define the composition \(\bigcap_{i\in I}x_{i}\) in the two following cases:

\begin{enumerate}
\def\labelenumi{(\alph{enumi})}
\item
  the set I is totally ordered;
\item
  E is associative and the \(x_{i}\) are pairwise permutable.
\end{enumerate}

In these two cases let S be the support of the family \((x_{i})\) . If J is a finite subset of I containing S, then \(\bigcap_{i\in J}x_{i} = \bigcap_{i\in S}x_{i}\) , as is seen by induction on the number of elements in J, applying Theorem 1 of §1 in case (a) and Theorem 3 of §1 in case (b). Let \(\bigcap_{i\in I}x_{i}\) denote the common value of the compositions \(\bigcap_{i\in J}x_{i}\) for all finite subsets of I containing S. When I is the interval \(\{p,\rightarrow\{\text{of }N\}\) , we also write \(\bigcap_{i=p}^{\infty}x_{i}\) .

With these definitions and notation, Theorems 1 and 3 of § 1 and the remarks following Theorem 3 extend to families with finite support.

The identity element of a law written additively is usually denoted by 0 and called zero or the null element (or sometimes the origin). Under a law written multiplicatively it is usually denoted by 1 and called the unit element (or unit).

\subsection{CANCELLABLE ELEMENTS}

DEFINITION 4. Given a law of composition \(\top\) on a set E, the mapping \(x \mapsto a \top x\) (resp. \(x \mapsto x \top a\) ) of E into itself is called left translation (resp. right translation) by an element \(a \in E\) .

On passing to the opposite law, left translations become right translations and conversely.

Let \(\pmb{\gamma}_{a},\pmb{\delta}_{a}\) (or \(\pmb {\gamma}(a),\pmb {\delta}(a)\) ) denote the left and right translations by \(a\in \mathbf{E}\) ; then

\[
\boldsymbol {\gamma} _ {a} (x) = a \top x, \quad \boldsymbol {\delta} _ {a} (x) = x \top a.
\]

PROPOSITION 1. If the law \(\top\) is associative, then for all \(x \in E\) and \(y \in E\)

\[
\gamma_ {x \top y} = \gamma_ {x} \circ \gamma_ {y}, \quad \delta_ {x \top y} = \delta_ {y} \circ \delta_ {x}.
\]

For all \(z \in \mathbf{E}\) :

\[
\boldsymbol {\Upsilon} _ {x \top y} (z) = (x \top y) \top z = x \top (y \top z) = \boldsymbol {\Upsilon} _ {x} (\boldsymbol {\Upsilon} _ {y} (z))
\]

\[
\boldsymbol {\delta} _ {x \top y} (z) = z \top (x \top y) = (z \top x) \top y = \boldsymbol {\delta} _ {y} (\boldsymbol {\delta} _ {x} (z))
\]

In other words, the mapping \(x \mapsto \gamma_x\) is a homomorphism from the magma E to the set \(\mathbf{E}^{\mathbb{E}}\) of mappings of E into itself with the law \((f, g) \mapsto f \circ g\) ; the mapping \(x \mapsto \delta_x\) is a homomorphism of E into the set \(\mathbf{E}^{\mathbb{E}}\) with the opposite law. If E is a monoid, these homomorphisms are unital.

DEFINITION 5. An element a of a magma E is called left (resp. right) cancellable (or regular) if left (resp. right) translation by a is injective. A left and right cancellable element is called a cancellable (or regular) element.

In other words, for \(a\) to be cancellable under the law \(\top\) , it is necessary and sufficient that each of the relations \(a \top x = a \top y\) , \(x \top a = y \top a\) imply \(x = y\) (it is said that ``a can be cancelled'' from each of these equalities). If there exists an identity element \(e\) under the law \(\top\) , it is cancellable under this law: the translations \(\gamma_e\) and \(\delta_e\) are then the identity mapping of E onto itself.

Examples. (1) Every natural number is cancellable under addition; every natural number \(\neq 0\) is cancellable under multiplication.

\begin{enumerate}
\def\labelenumi{(\arabic{enumi})}
\setcounter{enumi}{1}
\tightlist
\item
  In a lattice there can be no cancellable element under the law sup other than the identity element (least element) if it exists; similarly for inf. In particular, in the set of subsets of a set E, \(\varnothing\) is the only cancellable element under the law \(\cup\) and E the only cancellable element under the law \(\cap\) .
\end{enumerate}

PROPOSITION 2. The set of cancellable (resp. left cancellable, resp. right cancellable) elements of an associative magma is a submagma.

If \(\gamma_{x}\) and \(\gamma_{y}\) are injective so is \(\gamma_{x\top y} = \gamma_{x}\circ \gamma_{y}\) (Proposition 1). Similarly for \(\delta_{x\top y}\) .

\subsection{INVERTIBLE ELEMENTS}

DEFINITION 6. Let E be a unital magma, \(\top\) its law of composition, \(e\) its identity element and \(x\) and \(x'\) two elements of E. \(x'\) is called a left inverse (resp. right inverse, resp. inverse) of \(x\) if \(x' \top x = e\) (resp. \(x \top x' = e\) , resp. \(x' \top x = x \top x' = e\) ).

An element x of E is called left invertible (resp. right invertible, resp. invertible) if it has a left inverse (resp. right inverse, resp. inverse).

A monoid all of whose elements are invertible is called a group.

Symmetric and symmetrizable are sometimes used instead of inverse and invertible. When the law on E is written additively, we generally say negative instead of inverse.

Examples. (1) An identity element is its own inverse.

\begin{enumerate}
\def\labelenumi{(\arabic{enumi})}
\setcounter{enumi}{1}
\tightlist
\item
  In the set of mappings of E into E an element \(f\) is left invertible (resp. right invertible) if \(f\) is a surjection (resp. injection). The left inverses (resp. right inverses) are then the retractions (resp. sections) associated with \(f\) (Set Theory, II, § 3, no. 8, Definition 11). For \(f\) to be invertible it is necessary and sufficient that \(f\) be a bijection. Its unique inverse is then the inverse bijection of \(f\) .
\end{enumerate}

Let E and F be two unital magmas and f a unital homomorphism of E into F. If \(x'\) is the inverse of x in E, \(f(x')\) is the inverse of \(f(x)\) in F. Hence, if x is an invertible element of E, \(f(x)\) is an invertible element of F.

In particular, if R is an equivalence relation compatible with the law on a unital magma E, the canonical image of an invertible element of E is invertible in E/R.

PROPOSITION 3. Let E be a monoid and x an element of E.

\begin{enumerate}
\def\labelenumi{(\roman{enumi})}
\item
  For \(x\) to be left (resp. right) invertible it is necessary and sufficient that the right (resp. left) translation by \(x\) be surjective.
\item
  For \(x\) to be invertible it is necessary and sufficient that it be left and right invertible. In that case \(x\) has a unique inverse, which is also its unique left (resp. right) inverse.
\end{enumerate}

If \(x'\) is a left inverse of \(x\) then (Proposition 1)

\[
\boldsymbol {\delta} _ {x} \circ \boldsymbol {\delta} _ {x ^ {\prime}} = \boldsymbol {\delta} _ {x ^ {\prime} \top x} = \boldsymbol {\delta} _ {e} = \mathrm{Id} _ {\mathrm{E}}
\]

and \(\pmb{\delta}_{x}\) is surjective. Conversely, if \(\pmb{\delta}_{x}\) is surjective, there exists an element \(x^{\prime}\) of E such that \(\pmb{\delta}_{x}(x^{\prime}) = e\) and \(x^{\prime}\) is the left inverse of \(x\) . The other assertions of (i) follow similarly.

If \(x'\) (resp. \(x''\) ) is a left (resp. right) inverse of \(x\) , then

\[
x ^ {\prime} = x ^ {\prime} \top e = x ^ {\prime} \top (x \top x ^ {\prime \prime}) = (x ^ {\prime} \top x) \top x ^ {\prime \prime} = e \top x ^ {\prime \prime} = x ^ {\prime \prime},
\]

whence (ii).

Remark. Let E be a monoid and x an element of E. If x is left invertible it is left cancellable; for, if \(x'\) is a left inverse of x, then

\[
\gamma_ {x ^ {\prime}} \circ \gamma_ {x} = \gamma_ {x ^ {\prime} \top x} = \gamma_ {e} = \mathrm{Id} _ {\mathrm{E}}
\]

and \(\pmb{\gamma}_{x}\) is injective. In particular, if \(x\) is left invertible, the left and right translations by \(x\) are bijective. Conversely, suppose that \(\pmb{\gamma}_{x}\) is bijective; there exists \(x' \in E\) such that \(xx' = \pmb{\gamma}_{x}(x') = e\) ; then \(\pmb{\gamma}_{x}(x' x) = (xx') x = x = \pmb{\gamma}_{x}(e)\) and hence \(x' x = e\) , so that \(x\) is invertible. We see similarly that if \(\delta_{x}\) is bijective, \(x\) is invertible.

PROPOSITION 4. Let E be a monoid and x and y invertible elements of E with inverses \(x'\) and \(y'\) respectively. Then \(y' \top x'\) is the inverse of \(x \top y\) .

This follows from the relation

\[
\left(y ^ {\prime} \top x ^ {\prime}\right) \top (x \top y) = y ^ {\prime} \top \left(x ^ {\prime} \top x\right) \top y = y ^ {\prime} \top y = e
\]

and the analogous calculations for \((x\top y)\top (y^{\prime}\top x^{\prime})\)

COROLLARY 1. Let E be a monoid; if each of the elements \(x_{\alpha}\) of an ordered sequence \((x_{\alpha})_{\alpha \in A}\) of elements of E has an inverse \(x_{\alpha}'\) , the composition \(\bigcap_{\alpha \in A} x_{\alpha}\) has inverse \(\bigcap_{\alpha \in A'} x_{\alpha}'\) , where \(A'\) is the totally ordered set derived from A by replacing the order on A by the opposite order.

This corollary follows from Proposition 4 by induction on the number of elements in A.

In particular, if \(x\) and \(x'\) are inverses, \(\frac{\pi}{\top} x\) and \(\frac{\pi}{\top} x'\) are inverses for every integer \(n \geqslant 0\) .

COROLLARY 2. In a monoid the set of invertible elements is stable.

PROPOSITION 5. If in a monoid \(x\) and \(x'\) are inverses and \(x\) commutes with \(y\) , so does \(x'\) .

From \(x \top y = y \top x\) , we deduce \(x' \top (x \top y) \top x' = x' \top (y \top x) \top x'\) and hence \((x' \top x) \top (y \top x') = (x' \top y) \top (x \top x')\) , that is \(y \top x' = x' \top y\) .

COROLLARY 1. Let E be a monoid, X a subset of E and X' the centralizer of X. The inverse of every invertible element of X' belongs to X'.

COROLLARY 2. In a monoid the inverse of a central invertible element is a central element.

\subsection{MONOID OF FRACTIONS OF A COMMUTATIVE MONOID}

In this no., e will denote the identity element of a monoid and \(x^{*}\) the inverse of an invertible element x of E.

Let E be a commutative monoid, S a subset of E and S' the submonoid of E generated by S.

Lemma 1. In \(\mathbf{E} \times \mathbf{S}'\) the relation \(\mathbf{R}_{\xi}^{x}, y_{\xi}^{x}\) defined by:

``there exist \(a, b \in E\) and \(p, q, s \in S'\) such that \(x = (a, p)\) , \(y = (b, q)\) and \(aqs = bps\) ''

is an equivalence relation compatible with the law on the product monoid \(\mathbf{E} \times \mathbf{S}'\) .

It is immediate that R is reflexive and symmetric. Let \(x = (a, p)\) , \(y = (b, q)\) and \(z = (c, r)\) be elements of \(E \times S'\) such that \(R\xi x, y\xi\) and \(R\xi y, z\xi\) hold. Then there exist two elements s and t of \(S'\) such that

\[
a q s = b p s, \quad b r t = c q t,
\]

whence it follows that

\[
a r (s t q) = b p s r t = c p (s t q)
\]

and hence \(\mathbf{R}\{x, z\}\) holds, for \(stq\) belongs to \(\mathbf{S}'\) . The relation \(\mathbf{R}\) is therefore transitive.

Further, let \(x = (a, p), y = (b, q), x' = (a', p')\) , and \(y' = (b', q')\) be elements of \(E \times S'\) such that \(R\xi x, y\xi\) and \(R\xi x', y\xi\) hold. There exist \(s\) and \(s'\) in \(S'\) such that

\[
a q s = b p s, \quad a ^ {\prime} q ^ {\prime} s ^ {\prime} = b ^ {\prime} p ^ {\prime} s ^ {\prime},
\]

whence it follows that \((aa')(qq')(ss') = (bb')(pp')(ss')\) and hence \(\mathbf{R}\{\chi x', yy'\}\) for \(ss' \in S'\) . The equivalence relation \(\mathbf{R}\) is therefore compatible with the law of composition on \(\mathbf{E} \times \mathbf{S}'\) .

The quotient magma \((\mathbf{E} \times \mathbf{S}')/\mathbf{R}\) is a commutative monoid.

DEFINITION 7. Let E be a commutative monoid, S a subset of E and S' the submonoid of E generated by S. The quotient monoid (E × S')/R, where the equivalence relation R is as described in Lemma 1, is denoted by E \(_{S}\) and is called the monoid of fractions† of E associated with S (or with denominators in S).

For \(a \in \mathbf{E}\) and \(p \in \mathbf{S}'\) the class of \((a, p)\) modulo \(\mathbf{R}\) is in general denoted by \(a/p\) and called the fraction with numerator \(a\) and denominator \(p\) . Then by definition \((a/p)\) . \((a'/p') = aa'/pp'\) . The fractions \(a/p\) and \(a'/p'\) are equal if and only if there exists \(s\) in \(\mathbf{S}'\) with \(spa' = sp'a\) ; if so, there exist \(\sigma\) and \(\sigma'\) in \(\mathbf{S}'\) with \(a\sigma = a'\sigma'\) and \(p\sigma = p'\sigma'\) . In particular, \(a/p = sa/sp\) for \(a \in \mathbf{E}\) and \(s, p\) in \(\mathbf{S}'\) . The identity element of \(E_S\) is the fraction \(e/e\) .

Let \(a / e = \varepsilon(a)\) for all \(a \in E\) . The above shows that \(\varepsilon\) is a homomorphism of \(E\) into \(E_S\) , called canonical. For all \(p \in S'\) , \((p / e) \cdot (e / p) = e / e\) and hence \(e / p\) is the inverse of \(\varepsilon(p) = p / e\) ; every element of \(\varepsilon(S')\) is therefore invertible. Then \(a / p = (a / e) \cdot (e / p)\) , whence

\[
a / p = \varepsilon (a) \varepsilon (p) ^ {*}\tag{1}
\]

for \(a \in \mathbf{E}\) and \(p \in \mathbf{S}'\) , the monoid \(\mathbf{E}_{\mathbf{S}}\) is therefore generated by \(\varepsilon(\mathbf{E}) \cup \varepsilon(\mathbf{S})^*\) .

PROPOSITION 6. The notation is that of Definition 7 and \(\varepsilon\) denotes the canonical homomorphism of E into \(\mathbf{E}_{\mathrm{S}}\) .

\begin{enumerate}
\def\labelenumi{(\roman{enumi})}
\item
  Let \(a\) and \(b\) be in \(\mathbf{E}\) ; in order that \(\varepsilon(a) = \varepsilon(b)\) , it is necessary and sufficient that there exist \(s \in \mathbf{S}'\) with \(sa = sb\) .
\item
  For \(\varepsilon\) to be injective it is necessary and sufficient that every element of S be cancellable.
\item
  For \(\varepsilon\) to be bijective it is necessary and sufficient that every element of \(S\) be invertible.
\end{enumerate}

Assertion (i) is clear and shows that \(\varepsilon\) is injective if and only if every element of \(S'\) is cancellable; but since the set of cancellable elements of \(E\) is a submonoid of \(E\) (no. 2, Proposition 2), it amounts to the same to say that every element of \(S\) is cancellable.

If \(\varepsilon\) is bijective every element of S is invertible, for \(\varepsilon(S)\) is composed of invertible elements of \(E_{S}\) . Conversely, suppose every element of S is invertible; then every element of \(S'\) is invertible (no. 3, Corollary 2 to Proposition 4) and hence cancellable. Then \(\varepsilon\) is injective by (ii) and \(a/p = \varepsilon(a.p^{*})\) by (1), hence \(\varepsilon\) is surjective.

THEOREM 1. Let E be a commutative monoid, S a subset of E, \(E_{S}\) the monoid of fractions associated with S and \(\varepsilon: E \to E_{S}\) the canonical homomorphism. Further let f be a homomorphism of E into a monoid F (not necessarily commutative) such that every element of \(f(S)\) is invertible in F. There exists one and only one homomorphism \(\bar{f}\) of \(E_{S}\) into F such that \(f = \bar{f} \circ \varepsilon\) .

If \(\bar{f}\) is a homomorphism of \(\mathbf{E}_{\mathbf{S}}\) into \(\mathbf{F}\) such that \(f = \bar{f} \circ \varepsilon\) , then

\[
\bar {f} (a | p) = \bar {f} (\varepsilon (a) \varepsilon (p) ^ {*}) = \bar {f} (\varepsilon (a)) \bar {f} (\varepsilon (p)) ^ {*} = f (a) f (p) ^ {*}
\]

for \(a \in \mathbf{E}\) and \(p \in \mathbf{S}'\) , whence the uniqueness of \(\bar{f}\) .

Let \(g\) be the mapping of \(\mathbf{E} \times \mathbf{S}'\) into \(\mathbf{F}\) defined by \(g(a, p) = f(a).f(p)^*\) . We show that \(g\) is a homomorphism of \(\mathbf{E} \times \mathbf{S}'\) into \(\mathbf{F}\) . First of all,

\[
g (e, e) = f (e) f (e) ^ {*} = e.
\]

Let \((a, p)\) and \((a', p')\) be two elements of \(\mathbf{E} \times \mathbf{S}'\) ; as \(a'\) and \(p\) commute in \(\mathbf{E}\) , \(f(a')\) and \(f(p)\) commute in \(\mathbf{F}\) , whence \(f(a')f(p)^* = f(p)^*f(a')\) by no. 3, Proposition 5. Moreover \(f(p p')^* = f(p' p)^* = (f(p')f(p))^* = f(p)^*f(p')^*\) by no. 3, Proposition 4, whence

\[
\begin{array}{l} g (a a ^ {\prime}, p p ^ {\prime}) = f (a a ^ {\prime}) f (p p ^ {\prime}) ^ {*} = f (a) f (a ^ {\prime}) f (p) ^ {*} f (p ^ {\prime}) ^ {*} = f (a) f (p) ^ {*} f (a ^ {\prime}) f (p ^ {\prime}) ^ {*} \\ = g (a, p) g (a ^ {\prime}, p ^ {\prime}). \end{array}
\]

We show that \(g\) is compatible with the equivalence relation \(\mathbf{R}\) on \(\mathbf{E} \times \mathbf{S}'\) : if \((a, p)\) and \((a', p')\) are congruent mod. \(\mathbf{R}\) , there exists \(s \in \mathbf{S}'\) with \(spa' = sap'\) , whence \(f(s)f(p)f(a') = f(s)f(a)f(p')\) . As \(f(s)\) is invertible, it follows that \(f(p)f(a') = f(a)f(p')\) and then by left multiplication by \(f(p)^*\) and right multiplication by \(f(p')^*\)

\[
g (a ^ {\prime}, p ^ {\prime}) = f (a ^ {\prime}) f (p ^ {\prime}) ^ {*} = f (p) ^ {*} f (a) = f (a) f (p) ^ {*} = g (a, p).
\]

Hence there exists a homomorphism \(\bar{f}\) of \(\mathbf{E}_{\mathbf{S}}\) into \(\mathbf{F}\) such that \(\bar{f}(a/p) = g(a,p)\) , whence \(\bar{f}(\varepsilon(a)) = \bar{f}(a/e) = f(a)f(e)^* = f(a)\) . Hence \(\bar{f} \circ \varepsilon = f\) .

COROLLARY. Let E and F be two commutative monoids, S and T subsets of E and F respectively, f a homomorphism of E into F such that \(f(\mathrm{S}) \subset \mathrm{T}\) and \(\varepsilon: E \to E_{S}\) , \(\eta: F \to F_{T}\) the canonical homomorphisms. There exists one and only one homomorphism \(g: E_{S} \to F_{T}\) such that \(g \circ \varepsilon = \eta \circ f\) .

The homomorphism \(\eta \circ f\) of \(E\) into \(F_{T}\) maps every element of \(S\) to an invertible element of \(F_{T}\) .

Remarks. (1) Theorem 1 can also be expressed by saying that \((\mathbf{E}_{\mathbf{S}},\varepsilon)\) is the solution of the universal mapping problem for \(\mathbf{E}\) , relative to monoids, monoid homomorphisms and homomorphisms of \(\mathbf{E}\) into monoids which map the elements of \(\mathbf{S}\) to invertible elements (Set Theory, IV, § 3, no. 1). It follows (loc. cit.) that every other solution of this problem is isomorphic in a unique way to \((\mathbf{E}_{\mathbf{S}},\varepsilon)\) .

\begin{enumerate}
\def\labelenumi{(\arabic{enumi})}
\setcounter{enumi}{1}
\tightlist
\item
  For the existence of a solution to the above universal mapping problem it is unnecessary to assume that the monoid E is commutative, as follows from Set Theory, IV, § 3, no. 2 (cf.~Exercise 17).
\end{enumerate}

We mention two important special cases of monoids of fractions.

\begin{enumerate}
\def\labelenumi{(\alph{enumi})}
\item
  Let \(\overline{\mathbf{E}} = \mathbf{E}_{\mathbb{E}}\) . As the monoid \(\overline{\mathbf{E}}\) is generated by the set \(\varepsilon(\mathbf{E}) \cup \varepsilon(\mathbf{E})^*\) , which is composed of invertible elements, every element of \(\overline{\mathbf{E}}\) is invertible (no. 3, Corollary 2 to Proposition 4). In other words, \(\overline{\mathbf{E}}\) is a commutative group. Moreover, by Theorem 1 every homomorphism \(f\) of \(\mathbf{E}\) into a group \(\mathbf{G}\) can be uniquely factorized in the form \(f = \bar{f} \circ \varepsilon\) , where \(\bar{f}: \overline{\mathbf{E}} \to \mathbf{G}\) is a homomorphism. \(\overline{\mathbf{E}}\) is called the group of fractions of \(\mathbf{E}\) (or group of differences of \(\mathbf{E}\) in the case of additive notation).
\item
  Let \(\Phi = \mathrm{E}_{\Sigma}\) , where \(\Sigma\) consists of the cancellable elements of \(\mathbf{E}\) . By Proposition 6, (ii), the canonical homomorphism of \(\mathbf{E}\) into \(\Phi\) is injective; it will be profitable to identify \(\mathbf{E}\) with its image in \(\Phi\) . Hence \(\mathbf{E}\) is a submonoid of \(\Phi\) , every cancellable element of \(\mathbf{E}\) has an inverse in \(\Phi\) and every element of \(\Phi\) is of the form \(a/p = a.p^{*}\) with \(a \in \mathbf{E}\) and \(p \in \Sigma\) ; then \(a/p = a'/p'\) if and only if \(ap' = pa'\) . It is easily seen that the invertible elements of \(\Phi\) are the fractions \(a/p\) , where \(a\) and \(p\) are cancellable, and \(p/a\) is the inverse of \(a/p\) .
\end{enumerate}

Now let S be a set of cancellable elements of E and S' be the submonoid of E generated by S. If \(a/p\) and \(a'/p'\) are two elements of \(\mathbf{E}_{\mathbf{S}}\) , then \(a/p = a'/p'\) if and only if \(ap' = pa'\) (for \(sap' = spa'\) implies \(ap' = pa'\) for all \(s \in \mathbf{S}'\) ). \(\mathbf{E}_{\mathbf{S}}\) may therefore be identified with the submonoid of \(\Phi\) generated by \(\mathbf{E} \cup \mathbf{S}^*\) .

If every element of E is cancellable, then \(\Phi = \bar{E}\) and E is a submonoid of the commutative group \(\Phi\) . Conversely, if E is isomorphic to a submonoid of a group, every element of E is cancellable.

\subsection{APPLICATIONS: I. RATIONAL INTEGERS}

Consider the commutative monoid \(\mathbf{N}\) of natural numbers with law of composition addition; all the elements of \(\mathbf{N}\) are cancellable under this law (Set Theory, III, § 5, no. 2, Corollary 3). The group of differences of \(\mathbf{N}\) is denoted by \(\mathbf{Z}\) ; its elements are called the rational integers; its law is called addition of rational integers and also denoted by \(+\) . The canonical homomorphism from \(\mathbf{N}\) to \(\mathbf{Z}\) is injective and we shall identify each element of \(\mathbf{N}\) with its image in \(\mathbf{Z}\) . The elements of \(\mathbf{Z}\) are by definition the equivalence classes determined in \(\mathbf{N} \times \mathbf{N}\) by the relation between \((m_1, n_1)\) and \((m_2, n_2)\) which is written \(m_1 + n_2 = m_2 + n_1\) ; an element \(m\) of \(\mathbf{N}\) is identified with the class consisting of the elements \((m + n, n)\) , where \(n \in N\) ; it admits as negative in Z the class of elements \((n, m + n)\) . Every element \((p, q)\) of \(N \times N\) may be written in the form \((m + n, n)\) if \(p \geqslant q\) or in the form \((n, m + n)\) if \(p \leqslant q\) ; it follows that Z is the union of N and the set of negatives of the elements of N. The identity element 0 is the only element of N whose negative belongs to N.

For every natural number m, -m denotes the negative rational integer of m and -N denotes the set of elements -m for \(m \in N\) . Then

\[
\mathbf {Z} = \mathbf {N} \cup (- \mathbf {N}) \quad \text { and } \quad \mathbf {N} \cap (- \mathbf {N}) = \{0 \}.
\]

for \(m \in \mathbf{N}\) , \(m = -m\) if and only if \(m = 0\) .

Let \(m\) and \(n\) be two natural numbers;

\begin{enumerate}
\def\labelenumi{(\alph{enumi})}
\item
  if \(m \geqslant n\) , then \(m + (-n) = p\) , where \(p\) is the element of \(\mathbf{N}\) such that \(m = n + p\) ;
\item
  if \(m \leqslant n\) , then \(m + (-n) = -p\) , where \(p\) is the element of \(\mathbf{N}\) such that \(m + p = n\) ;
\end{enumerate}

\[
(c) (- m) + (- n) = - (m + n).
\]

Properties (b) and (c) follow from no. 3, Proposition 4; as \(\mathbf{Z} = \mathbf{N} \cup (-\mathbf{N})\) , addition in \(\mathbf{N}\) and properties (a), (b) and (c) describe completely addition in \(\mathbf{Z}\) .

More generally -x is used to denote the negative of an arbitrary rational integer x; the composition \(x + (-y)\) is abbreviated to x - y (cf.~no. 8).

The order relation \(\leqslant\) between natural numbers is characterized by the following property: \(m \leqslant n\) if and only if there exists an integer \(p \in N\) such that \(m + p = n\) (Set Theory, III, §3, no. 6, Proposition 13 and §5, no. 2, Proposition 2). The relation \(y - x \in N\) between rational integers x and y is a total order relation on Z which extends the order relation \(\leqslant\) on N. For, for all \(x \in Z\) , \(x - x = 0 \in N\) ; if \(y - x \in N\) and \(z - y \in N\) , then

\[
z - x = (z - y) + (y - x) \in \mathbf {N},
\]

for \(\mathbf{N}\) is stable under addition; if \(y - x \in \mathbf{N}\) and \(x - y \in \mathbf{N}\) , then \(y - x = 0\) , for \(0\) is the only element of \(\mathbf{N}\) whose negative belongs to \(\mathbf{N}\) ; for arbitrary rational integers \(x\) and \(y, y - x \in \mathbf{N}\) or \(x - y \in \mathbf{N}\) , for \(\mathbf{Z} = \mathbf{N} \cup (-\mathbf{N})\) ; finally, if \(x\) and \(y\) are natural numbers, then \(y - x \in \mathbf{N}\) if and only if there exists \(p \in \mathbf{N}\) such that \(x + p = y\) . This order relation is also denoted by \(\leqslant\) .

Henceforth when Z is considered as an ordered set, it will always be, unless otherwise mentioned, with the ordering that has just been defined, the natural numbers are identified with the integers \(\geqslant 0\) ; they are also called positive integers; the integers \(\leqslant 0\) , negatives of the positive integers, are called negative integers; the integers \textgreater{} 0 (resp. \textless{} 0) are called strictly positive (resp. strictly negative); the set of integers \textgreater{} 0 is sometimes denoted by N*.

Let \(x, y\) and \(z\) be three rational integers; then \(x \leqslant y\) if and only if \(x + z \leqslant y + z\) . For \(x - y = (x + z) - (y + z)\) . This property is expressed by saying that the order relation on Z is invariant under translation.

\subsection{APPLICATIONS: II. MULTIPLICATION OF RATIONAL INTEGERS}

Lemma 2. Let \(\mathbf{E}\) be a monoid and \(x\) an element of \(\mathbf{E}\) .

\begin{enumerate}
\def\labelenumi{(\roman{enumi})}
\item
  There exists a unique homomorphism \(f\) of \(\mathbf{N}\) into \(\mathbf{E}\) with \(f(1) = x\) and \(f(n) = \frac{n}{T} x\) for all \(n \in \mathbf{N}\) .
\item
  If \(x\) is invertible, there exists a unique homomorphism \(g\) of \(\mathbf{Z}\) into \(\mathbf{E}\) such that \(g(1) = x\) and \(g\) coincides with \(f\) on \(\mathbf{N}\) .
\end{enumerate}

Writing \(f(n) = \top^n x\) for all \(n \in \mathbf{N}\) , the formulae

\[
\overset {0} {\top} x = e \quad \text { and } \quad \left(\overset {m} {\top} x\right) \top \left(\overset {n} {\top} x\right) = \overset {m + n} {\top} x
\]

(no. 1) express the fact that \(f\) is a homomorphism of \(\mathbf{N}\) into \(\mathbf{E}\) and obviously \(f(1) = x\) . If \(f'\) is a homomorphism of \(\mathbf{N}\) into \(\mathbf{E}\) such that \(f'(1) = x\) , then \(f = f'\) , by § 1, no. 4, Proposition 1, (iv).

Suppose now that \(x\) is invertible. By no. 3, Corollary 2 to Proposition 4, \(f(n) = \overset{n}{\top} x\) is invertible for every integer \(n \geqslant 0\) . By construction, \(\mathbf{Z}\) is the group of differences of \(\mathbf{N}\) and hence (no. 4, Theorem 1) \(f\) extends uniquely to a homomorphism \(g\) of \(\mathbf{Z}\) into \(\mathbf{E}\) . If \(g'\) is a homomorphism of \(\mathbf{Z}\) into \(\mathbf{E}\) with \(g'(1) = x\) , the restriction \(f'\) of \(g'\) to \(\mathbf{N}\) is a homomorphism of \(\mathbf{N}\) into \(\mathbf{E}\) with \(f'(1) = x\) . Hence \(f' = f\) , whence \(g' = g\) .

We shall apply Lemma 2 to the case where the monoid E is Z; for every integer \(m \in Z\) there therefore exists an endomorphism \(f_{m}\) of Z characterized by \(f_{m}(1) = m\) . If m is in N, the mapping \(n \mapsto mn\) of N into N is an endomorphism of the magma N (Set Theory, III, §3, no. 3, Corollary to Proposition 5); hence \(f_{m}(n) = mn\) for all m, n in N.

Multiplication on \(\mathbf{N}\) can therefore be extended to multiplication on \(\mathbf{Z}\) by the formula \(mn = f_m(n)\) for \(m, n\) in \(\mathbf{Z}\) . We shall establish the formulae:

\[
x y = y x\tag{2}
\]

\[
(x y) z = x (y z)\tag{3}
\]

\[
x (y + z) = x y + x z\tag{4}
\]

\[
(x + y) z = x z + y z\tag{5}
\]

\[
0. x = x. 0 = 0\tag{6}
\]

\[
1. x = x. 1 = x\tag{7}
\]

\[
(- 1). x = x. (- 1) = - x\tag{8}
\]

for \(x, y, z\) in \(\mathbf{Z}\) . (*In other words, \(\mathbf{Z}\) is a commutative ring.*) The formulae \(x(y + z) = xy + xz\) and \(x.0 = 0\) express the fact that \(f_x\) is an endomorphism of the additive monoid \(\mathbf{Z}\) and \(f_{x}(1) = x\) may be written \(x.1 = x\) . The endomorphism \(f_{x} \circ f_{y}\) of \(\mathbf{Z}\) maps 1 to \(xy\) and hence is equal to \(f_{xy}\) , whence (3). Now \(f_{x}(-y) = -f_{x}(y)\) , that is \(x(-y) = -xy\) ; similarly, the endomorphism \(y \mapsto -xy\) of \(\mathbf{Z}\) maps 1 to \(-x\) , whence \((-x).y = -xy\) and therefore

\[
(- x) (- y) = - (x (- y)) = - (- x y) = x y.
\]

For \(m, n\) in \(\mathbf{N}\) , \(mn = nm\) (Set Theory, III, § 3, no. 3, Corollary to Proposition 5), whence \((-m).n = n(-m)\) and \((-m)(-n) = (-n)(-m)\) ; as \(\mathbf{Z} = \mathbf{N} \cup (-\mathbf{N})\) , \(xy = yx\) for \(x, y\) in \(\mathbf{Z}\) ; and this formula means that (5) follows from (4) and completes the proof of formula (6) to (8).

\subsection{APPLICATIONS: III. GENERALIZED POWERS}

Let E be a monoid with identity element e and law of composition denoted by T. If x is invertible in E, let \(g_{x}\) be the homomorphism of Z into E mapping 1 to x. Let \(g_{x}(n) = \frac{n}{T} x\) for all \(n \in Z\) ; by Lemma 2 this notation is compatible for \(n \in N\) with the notation introduced earlier. Then

\[
{ } ^ { m + n } \top x = \left( \begin{array} { c } m \\ T \end{array} x \right) \top \left( \begin{array} { c } n \\ T \end{array} x \right)\tag{9}
\]

\[
\frac {0}{T} x = e\tag{10}
\]

\[
\frac {1}{T} x = x\tag{11}
\]

for \(x\) invertible in \(\mathbf{E}\) and \(m, n\) in \(\mathbf{Z}\) . Further, if \(y = \frac{m}{T} x\) , the mapping \(n \mapsto g_x(mn)\) of \(\mathbf{Z}\) into \(\mathbf{E}\) is a homomorphism mapping 1 to \(y\) , whence \(g_x(mn) = g_y(n)\) , that is

\[
\frac {m n}{T} x = \frac {n}{T} \left(\frac {m}{T} x\right).\tag{12}
\]

As \(-1\) is the negative of \(1\) in \(\mathbf{Z}\) , \(\overline{\top}^{-1}x\) is the inverse of \(x = \frac{1}{\top} x\) in E. If we write \(n = -m\) in (9), it is seen that \(\overline{\top}^{-m}x\) is the inverse of \(\overline{\top}^{m}x\) .

\subsection{NOTATION}

\begin{enumerate}
\def\labelenumi{(\alph{enumi})}
\tightlist
\item
  As a general rule the law of a commutative monoid is written additively. It is then a convention that -x denotes the negative of x. The notation \(x + (-y)\) is abbreviated to x - y and similarly
\end{enumerate}

\[
x + y - z, \quad x - y - z, \quad x - y + z - t, \quad \text { etc. } \dots
\]

represent respectively

\[
x + y + (- z), \quad x + (- y) + (- z), \quad x + (- y) + z + (- t), \quad \text { etc. } \dots
\]

For \(n \in \mathbf{Z}\) the notation \(\top^n x\) is replaced by \(nx\) . Formulae (9) to (12) then become

\[
(m + n). x = m. x + n. x\tag{13}
\]

\[
0. x = 0\tag{14}
\]

\[
1. x = x\tag{15}
\]

\[
m. (n. x) = (m n). x\tag{16}
\]

where m and n belong to N or even to Z if x admits a negative. Also in the latter case the relation \((-1)\) . x = -x holds. We also note the formula

\[
n. (x + y) = n. x + n. y.\tag{17}
\]

\begin{enumerate}
\def\labelenumi{(\alph{enumi})}
\setcounter{enumi}{1}
\tightlist
\item
  Let \(\mathbf{E}\) be a monoid written multiplicatively. For \(n \in \mathbf{Z}\) the notation \(\frac{n}{T} x\) is replaced by \(x^n\) . We have the relations
\end{enumerate}

\[
\begin{array}{r l} x ^ {m + n} & = x ^ {m}. x ^ {n} \\ x ^ {0} & = 1 \\ x ^ {1} & = x \\ (x ^ {m}) ^ {n} & = x ^ {m n} \end{array}
\]

and also \((xy)^n = x^n y^n\) if \(x\) and \(y\) commute.

When \(x\) has an inverse, this is precisely \(x^{-1}\) . The notation \(\frac{1}{x}\) is also used instead of \(x^{-1}\) . Finally, when the monoid \(E\) is commutative, \(\frac{x}{y}\) or \(x / y\) is also used for \(xy^{-1}\) .

\section{ACTIONS}

\subsection{ACTIONS}

DEFINITION 1. Let \(\Omega\) and \(\mathbf{E}\) be two sets. A mapping of \(\Omega\) into the set \(\mathbf{E}^{\mathbf{E}}\) of mappings of \(\mathbf{E}\) into itself is called an action of \(\Omega\) on \(\mathbf{E}\) .

Let \(\alpha \mapsto f_{\alpha}\) be an action of \(\Omega\) on E. The mapping \((\alpha, x) \mapsto f_{\alpha}(x)\) (resp. \((x, \alpha) \mapsto f_{\alpha}(x)\) ) is called the law of left (resp. right) action of \(\Omega\) on E† associated with the given action of \(\Omega\) on E. Given a mapping \(g\) of \(\Omega \times E\) (resp. E × \(\Omega\) ) into E, there exists one and only one action \(\alpha \mapsto f_{\alpha}\) of \(\Omega\) on E such that the associated law of left (resp. right) action is \(g\) (Set Theory, II, § 5, no. 2, Proposition 3).

In this chapter we shall say, for the sake of abbreviation, ``law of action''

instead of ``law of left action''. The element \(f_{\alpha}(x)\) of E (for \(\alpha \in \Omega\) and \(x \in E\) ) is sometimes called the transform of x under \(\alpha\) or the composition of \(\alpha\) and x. It is often denoted by left (resp. right) multiplicative notation \(\alpha.x\) (resp. \(x.\alpha\) ), the dot may be omitted; the composition of \(\alpha\) and x is then called the product of \(\alpha\) and x (resp. x and \(\alpha\) ). The exponential notation \(x^{\alpha}\) is also used. In the arguments of the following paragraphs we shall generally use the notation \(\alpha \perp x\) . The elements of \(\Omega\) are often called operators.

Examples. (1) Let E be an associative magma written multiplicatively. The mapping which associates with a strictly positive integer n the mapping \(x \mapsto x^{n}\) of E into itself is an action of \(N^{*}\) on E. If E is a group, the mapping which associates with a rational integer a the mapping \(x \mapsto x^{a}\) of E into E is an action of Z on E.

\begin{enumerate}
\def\labelenumi{(\arabic{enumi})}
\setcounter{enumi}{1}
\item
  Let \(\mathbf{E}\) be a magma with law denoted by \(\top\) . The mapping which associates with \(x \in \mathbf{E}\) the mapping \(\mathbf{A} \mapsto x \top \mathbf{A}\) of the set of subsets of \(\mathbf{E}\) into itself is an action of \(\mathbf{E}\) on \(\mathfrak{P}(\mathbf{E})\) .
\item
  Let E be a set. The identity mapping of \(E^{E}\) is an action of \(E^{E}\) on E, called the canonical action. The corresponding law of action is the mapping \((f, x) \mapsto f(x)\) of \(E^{E} \times E\) into E.
\item
  Let \((\Omega_{i})_{i\in I}\) be a family of sets. For all \(i\in I\) , let \(f_{i}\colon \Omega_{i}\to E^{\mathbb{E}}\) be an action of \(\Omega_{i}\) on E. Let \(\Omega\) be the sum of the \(\Omega_{i}\) (Set Theory, II, § 4, no. 8). The mapping \(f\) of \(\Omega\) onto \(E^{\mathbb{E}}\) , extending the \(f_{i}\) , is an action of \(\Omega\) on E. This allows us to reduce the study of a family of actions to that of a single action.
\item
  Given an action of \(\Omega\) on \(\mathbf{E}\) with law denoted by \(\bot\) , a subset \(\Xi\) of \(\Omega\) and a subset \(\mathbf{X}\) of \(\mathbf{E}\) , \(\Xi \perp \mathbf{X}\) denotes the set of \(\alpha \perp x\) with \(\alpha \in \Xi\) and \(x \in \mathbf{X}\) ; when \(\Xi\) consists of a single element \(\alpha\) , we generally write \(\alpha \perp \mathbf{X}\) instead of \(\{\alpha\} \perp \mathbf{X}\) . The mapping which associates with \(\alpha \in \Omega\) the mapping \(\mathbf{X} \mapsto \alpha \perp \mathbf{X}\) is an action of \(\Omega\) on \(\mathfrak{P}(\mathbf{E})\) , which is said to be derived from the given action by extension to the set of subsets.
\item
  Let \(\alpha \mapsto f_{\alpha}\) be an action of \(\Omega\) on \(\mathbf{E}\) . Let \(g\) be a mapping of \(\Omega'\) into \(\Omega\) . Then the mapping \(\beta \mapsto f_{g(\beta)}\) is an action of \(\Omega'\) on \(\mathbf{E}\) .
\item
  Let \(f: \mathbf{E} \times \mathbf{E} \to \mathbf{E}\) be a law of composition on a set \(\mathbf{E}\) . The mapping \(\gamma: x \mapsto \gamma_x\) (resp. \(\delta: x \mapsto \delta_x\) ) (§ 2, no. 2) which associates with the element \(x \in \mathbf{E}\) left (resp. right) translation by \(x\) is an action of \(\mathbf{E}\) on itself; it is called the left (resp. right) action of \(\mathbf{E}\) on itself derived from the given law. When \(f\) is commutative, these two actions coincide.
\end{enumerate}

The law of left (resp. right) action associated with \(\gamma\) is f (resp. the opposite law to f). The law of right (resp. left) action associated with \(\delta\) is f (resp. the opposite law to f).

Let \(\Omega, \mathbf{E}, \mathbf{F}\) be sets, \(\alpha \mapsto f_{\alpha}\) an action of \(\Omega\) on \(\mathbf{E}\) and \(\alpha \mapsto g_{\alpha}\) an action of \(\Omega\) on \(\mathbf{F}\) . An \(\Omega\) -morphism of \(\mathbf{E}\) into \(\mathbf{F}\) , or mapping of \(\mathbf{E}\) into \(\mathbf{F}\) compatible with the action of \(\Omega\) , is a mapping \(h\) of \(\mathbf{E}\) into \(\mathbf{F}\) such that

\[
g _ {\alpha} (h (x)) = h \left(f _ {\alpha} (x)\right)
\]

for all \(\alpha \in \Omega\) and \(x \in E\) . The composition of two \(\Omega\) -morphisms is an \(\Omega\) -morphism.

Let \(\Omega, \Xi, E, F\) be sets, \(\alpha \mapsto f_{\alpha}\) an action of \(\Omega\) on \(E\) , \(\beta \mapsto g_{\beta}\) an action of \(\Xi\) on \(F\) and \(\phi\) a mapping of \(\Omega\) into \(\Omega'\) . A \(\phi\) -morphism of \(E\) into \(F\) is a mapping \(h\) of \(E\) into \(F\) such that

\[
g _ {\phi (\alpha)} (h (x)) = h (f _ {\alpha} (x))
\]

for all \(\alpha \in \Omega\) and \(x \in \mathbf{E}\) .

\subsection{SUBSETS STABLE UNDER AN ACTION. INDUCED ACTION}

DEFINITION 2. A subset \(\mathbf{A}\) of a set \(\mathbf{E}\) is called stable under an action \(\alpha \mapsto f_{\alpha}\) of \(\Omega\) on \(\mathbf{E}\) if \(f_{\alpha}(\mathbf{A}) \subset \mathbf{A}\) for all \(\alpha \in \Omega\) . An element \(x\) of \(\mathbf{E}\) is called invariant under an element \(\alpha\) of \(\Omega\) if \(f_{\alpha}(x) = x\) .

The intersection of a family of stable subsets of E under a given action is stable. There therefore exists a smallest stable subset of E containing a given subset X of E; it is said to be generated by X; it consists of the elements \((f_{\alpha_{1}} \circ f_{\alpha_{2}} \circ \cdots \circ f_{\alpha_{n}})(x)\) , where \(x \in X, n \geqslant 0, \alpha_{1} \in \Omega\) for all i.

Remark. Let E be a magma with law denoted by \(\top\) . It should be noted that a subset A of E which is stable under the left action on E on itself is not necessarily stable under the right action of E on itself; a subset A of E stable under the left (resp. right) action of E on itself is stable under the law on E but the converse is not in general true. More precisely, A is stable under the law on E if and only if \(A \top A \subset A\) whereas A is stable under the left (resp. right) action on E on itself if and only if \(E \top A \subset A\) (resp. \(A \top E \subset A\) ).

Example. Take the magma E to be the set N with multiplication. The set \(\{1\}\) is stable under the internal law of N, but the stable subset under the action of N on itself generated by \(\{1\}\) is the whole of N.

DEFINITION 3. Let \(\alpha \mapsto f_{\alpha}\) be an action of \(\Omega\) on E and A a stable subset of E. The mapping which associates with an element \(\alpha \in \Omega\) the restriction of \(f_{\alpha}\) to A (considered as a mapping of A into itself) is an action of \(\Omega\) on A said to be induced by the given action.

\subsection{QUOTIENT ACTION}

DEFINITION 4. Let \(\alpha \mapsto f_{\alpha}\) be an action of a set \(\Omega\) on a set \(\mathbf{E}\) . An equivalence relation \(\mathbf{R}\) on \(\mathbf{E}\) is said to be compatible with the given action if, for all elements \(x\) and \(y\) of \(\mathbf{E}\) such that \(x \equiv y \pmod{\mathbf{R}}\) and all \(\alpha \in \Omega\) , \(f_{\alpha}(x) \equiv f_{\alpha}(y) \pmod{\mathbf{R}}\) . The mapping which associates with an element \(\alpha \in \Omega\) the mapping of \(\mathbf{E} / \mathbf{R}\) into itself derived from \(f_{\alpha}\) by passing to the quotients is an action of \(\Omega\) on \(\mathbf{E} / \mathbf{R}\) called the quotient of the action of \(\Omega\) on \(\mathbf{E}\) .

Let E be a magma and R an equivalence relation on E. R is said to be left (resp. right) compatible with the law on E if it is compatible with the left (resp.

(right) action of E on itself derived from the law on E. For R to be compatible with the law on E it is necessary and sufficient that it be left and right compatible with the law on E.

We leave to the reader the statement and proof of the analogues of Propositions 6, 7 and 8 of § 1, no. 6.
